\section{Discussion}
\label{sec:discussion}

The central result of this work is the throughput preservation under bursty traffic shown in Section~\ref{sec:results}-A: a 3.1$\times$ throughput improvement over a representative fixed-threshold baseline, at zero mode transitions versus nineteen. This directly addresses the thrashing failure mode identified in Section~\ref{sec:related} and confirms the central novelty claim relative to \cite{agarwal2006}.

This benefit is not free. Real post-route synthesis shows a 13.6\% area overhead for the additional comparator and threshold-widening logic -- below the originally targeted 10\% ``lightweight'' threshold, but a substantial improvement over the 20--26\% suggested by pre-layout estimation alone. The appropriate framing for this result is a bounded, quantified trade-off rather than a cost-free improvement: comparator-based adaptive control achieves a meaningful behavioral improvement at a modest, real silicon cost, without resorting to the prediction-hardware overhead of learned approaches such as \cite{karthikeyan2025}. The 104 additional standard cells required for the adaptive predictor logic represent a minimal hardware investment compared to the ML inference hardware required by \cite{karthikeyan2025}.

Compared to prior work: versus \cite{agarwal2006}, this work adds the activity-rate second condition and eliminates thrashing under burst workloads. Versus \cite{karthikeyan2025}, this work achieves comparable qualitative adaptivity -- suppressing sleep during bursty periods, entering sleep during genuinely idle periods -- with no inference hardware, no learned models, and no per-cycle software classification. Versus \cite{chang2012}, this work's contribution is orthogonal: Chang et al.'s two-phase wake-up mechanism is adopted unchanged in both baseline and proposed designs as a shared constant.

\noindent\textit{Limitations:}
This work's energy comparisons use a derived metric -- transitions weighted by maximum wake-up latency -- rather than a first-principles joule measurement; a SPICE-level or annotated-switching-activity power simulation would strengthen this claim further. Post-route results are reported at a single LP corner; multi-corner PVT analysis was outside this project's scope. The synthetic workloads, while designed to specifically probe the threshold-crossing behavior under test, are not drawn from a real application trace. DRC verification was limited by the available PDK configuration on the server environment; the 7,376 spacing violations reported by the router reflect missing DRC deck rules rather than functional errors.
